\documentclass[11pt,a4paper]{article}
\usepackage[UTF8]{ctex}
\usepackage{amsmath,amssymb}
\usepackage{booktabs}
\usepackage{hyperref}
\usepackage{geometry}
\usepackage{enumitem}
\usepackage{xcolor}
\usepackage{array}
\usepackage{multirow}

\geometry{margin=2.5cm}
\hypersetup{colorlinks=true,linkcolor=blue,urlcolor=blue,citecolor=blue}

\title{联合训练进展汇报}
\author{}
\date{2026-03-16}

\begin{document}
\maketitle

\section{框架选型}

调研了两个可行的 RL 框架：

\begin{table}[h]
\centering
\begin{tabular}{llll}
\toprule
\textbf{框架} & \textbf{训练后端} & \textbf{Rollout 引擎} & \textbf{状态} \\
\midrule
\textbf{verl} & FSDP (ZeRO-3) & vLLM & \textbf{已选用} \\
slime & Megatron & SGLang & 暂不使用 \\
\bottomrule
\end{tabular}
\end{table}

选择 verl 的原因：

% TODO: 补充选型理由

\section{已完成的工作}

\subsection{Infra：联合训练的 Logit 融合}

核心设计是在输出层进行 logit 线性组合：
\begin{equation}
\text{logits}_{\text{fused}} = (1 - \lambda) \cdot \text{logits}_{\text{model}_1} + \lambda \cdot \text{logits}_{\text{model}_2}
\end{equation}

其中 $\text{model}_1$ 为锚定模型（frozen），$\text{model}_2$ 为可训练模型。两个子模型独立前向传播，仅在 logit 层融合，没有 hidden states 的交互。融合后的 logits 用于计算 RL 损失并反向传播——因为 $\text{model}_1$ 的参数被冻结，梯度只更新 $\text{model}_2$。

实现要点：
\begin{itemize}[nosep]
    \item 模型类 \texttt{QwenJointForCausalLM}：支持 HuggingFace \texttt{from\_pretrained} 加载，兼容 FSDP 分片
    \item vLLM rollout 集成：训练时用融合 logits 采样，评估时提取 model2 独立推理
    \item 已完成 25 项语义测试（融合计算、old-log-prob 重算、权重提取等）
\end{itemize}

\subsection{损失函数：从 GRPO 切换到 MiniRL}

最初用标准 GRPO（vanilla 模式）跑联合训练，在 GSM8K 上 0\% 准确率（EXP-01），没有学到任何东西。切换到 \textbf{MiniRL}~\cite{minirl} 后，训练开始收敛。

\subsubsection{MiniRL 算法简介}

MiniRL 的核心改动是用\textbf{二值裁剪掩码}替代 PPO/GRPO 的 min-clip 机制。标准 GRPO 的 policy loss：
\begin{equation}
\mathcal{L}_{\text{GRPO}} = -\mathbb{E}\left[\min\left(r_t \hat{A}_t, \; \text{clip}(r_t, 1-\epsilon, 1+\epsilon) \hat{A}_t\right)\right]
\end{equation}
其中 $r_t = \frac{\pi_\theta(a_t|s_t)}{\pi_{\theta_{\text{old}}}(a_t|s_t)}$ 是重要性采样比率。

MiniRL 将其简化为：
\begin{equation}
\mathcal{L}_{\text{MiniRL}} = -\mathbb{E}\left[m_t \cdot \hat{A}_t \cdot \log \pi_\theta(a_t|s_t)\right]
\end{equation}

其中 $m_t$ 是二值掩码：
\begin{equation}
m_t = \begin{cases}
0, & \text{if } \hat{A}_t > 0 \text{ and } r_t > 1 + \epsilon_{\text{high}} \\
0, & \text{if } \hat{A}_t < 0 \text{ and } r_t < 1 - \epsilon_{\text{low}} \\
1, & \text{otherwise}
\end{cases}
\end{equation}

\textbf{与 GRPO 的关键区别}：

\begin{table}[h]
\centering
\begin{tabular}{p{3cm}p{5.5cm}p{5.5cm}}
\toprule
\textbf{特性} & \textbf{GRPO (vanilla)} & \textbf{MiniRL} \\
\midrule
裁剪方式 & $\min(r_t \hat{A}_t, \text{clip}(r_t) \hat{A}_t)$ & 二值掩码 $m_t \in \{0, 1\}$ \\
梯度通路 & 通过 ratio 和 log\_prob & 仅通过 $\log \pi_\theta$（REINFORCE 风格） \\
损失聚合 & token-mean（除以 token 数） & seq-mean-token-sum（不做 token 级归一化） \\
Advantage & 可以 attach 梯度 & 必须 detach，作为常数 \\
\bottomrule
\end{tabular}
\end{table}

MiniRL 本质上回到了 REINFORCE 的梯度形式，但加了一个基于 policy staleness 的二值门控。它的理论优势是在 \texttt{seq-mean-token-sum} 聚合下保持了梯度估计的无偏性。

\subsubsection{Advantage 计算：Dr.GRPO}

我们使用 Dr.GRPO 变体来计算 advantage。标准 GRPO 对每个 prompt 的 $G$ 条响应做归一化：
\begin{equation}
\hat{A}_i^{\text{GRPO}} = \frac{r_i - \text{mean}(\mathbf{r})}{\text{std}(\mathbf{r}) + \epsilon}
\end{equation}

Dr.GRPO 去掉了标准差归一化，直接用均值中心化：
\begin{equation}
\hat{A}_i^{\text{Dr.GRPO}} = r_i - \text{mean}(\mathbf{r})
\end{equation}
其中 $r_i \in \{-1, +1\}$ 是二值奖励。去掉标准差的动机是：当一组响应全对或全错时，$\text{std}(\mathbf{r}) = 0$，标准 GRPO 的归一化会产生数值不稳定；而且对于二值奖励，除以标准差不提供额外信息，反而引入不必要的缩放。

\subsection{实验结果}

在 H800 服务器（8 GPU）上用 Qwen3-1.7B-Base 训练，速度很快（100 步约 1.5 小时）。共跑了 6 组实验：

\subsubsection{实验总览}

\begin{table}[h]
\centering
\begin{tabular}{lllccc}
\toprule
\textbf{实验} & \textbf{模型} & \textbf{算法} & \textbf{联合训练} & \textbf{MATH-500 最佳} & \textbf{状态} \\
\midrule
EXP-01 & 1.7B & GRPO & 是 ($\lambda$=0.5) & 0\% & 失败 \\
EXP-02 & 4B & GRPO+IS & 是 ($\lambda$=0.5) & N/A & 完成 \\
EXP-03 & 4B & GRPO+IS & 是 ($\lambda$=0.5) & N/A & 中断 \\
EXP-04 & 1.7B & MiniRL & 是 ($\lambda$=0.55) & \textbf{63.0\%} & 完成 \\
EXP-05 & 1.7B & MiniRL & 否 & \textbf{64.0\%} & 完成 \\
EXP-06 & 1.7B & MiniRL & 否 & 待运行 & 计划中 \\
\bottomrule
\end{tabular}
\end{table}

\subsubsection{离线评估对比（EVAL-01 vs EVAL-02）}

用 vLLM 离线推理，对 EXP-04（联合训练）和 EXP-05（单模型 baseline）在 5 个数学基准上评测（$n=3$，$T=1.0$）：

\begin{table}[h]
\centering
\begin{tabular}{lrccc}
\toprule
\textbf{Benchmark} & \textbf{样本数} & \textbf{联合训练 mean@3} & \textbf{Baseline mean@3} & \textbf{差值} \\
\midrule
\textbf{MATH-500} & 500 & \textbf{64.2\%} & 61.4\% & \textbf{+2.8\%} \\
AIME-2025 & 30 & 4.4\% & 5.6\% & $-$1.2\% \\
AMC-2023 & 40 & 36.7\% & 41.7\% & $-$5.0\% \\
MinervaMAth & 272 & 24.4\% & 27.6\% & $-$3.2\% \\
OlympiadBench & 674 & 28.3\% & 28.1\% & +0.2\% \\
\bottomrule
\end{tabular}
\end{table}

\textbf{观察}：联合训练在训练域（MATH-500）上优于 baseline（+2.8\%），但在 OOD 基准上整体弱于 baseline。不过这个对比\textbf{不公平}——baseline 存在严重的梯度裁剪问题（见下文），实际上没有充分训练。

\subsection{发现的关键问题：梯度裁剪瓶颈}

在分析 EXP-05（baseline）的训练日志时，发现了一个严重的超参数不匹配：

\textbf{\texttt{grad\_clip=1.0} 对 MiniRL 来说太激进了。}

MiniRL 使用 \texttt{seq-mean-token-sum} 聚合（token 级损失求和，不除以序列长度），导致梯度范数天然比 GRPO 的 \texttt{token-mean} 大 100--500 倍：

\begin{table}[h]
\centering
\begin{tabular}{lccc}
\toprule
\textbf{聚合模式} & \textbf{典型 grad\_norm} & \textbf{裁剪率} & \textbf{有效学习率} \\
\midrule
token-mean (GRPO) & 0.5 -- 5 & $\sim$20--50\% & $\approx 10^{-6}$ \\
seq-mean-token-sum (MiniRL) & 231 -- 515 & \textbf{100\%} & $\approx 2.8 \times 10^{-9}$ \\
\bottomrule
\end{tabular}
\end{table}

这意味着 MiniRL 的有效学习率被压缩了约 300 倍：
\begin{equation}
\text{lr}_{\text{effective}} = \text{lr} \times \frac{\text{grad\_clip}}{\|\nabla\|} = 10^{-6} \times \frac{1.0}{368} \approx 2.8 \times 10^{-9}
\end{equation}

这解释了 MATH-500 准确率在 step 40 之后就进入平台期（$\sim$60\%），后续 125 步只涨了约 4 个百分点。

\textbf{修复方案}：将 \texttt{grad\_clip} 从 1.0 提高到 500.0。在已观测到的 33 个 step 中，只有 1 个 step 的 grad\_norm 超过 500（step 5 的 515.5），其余都在 500 以内。EXP-06 就是用 \texttt{grad\_clip=500.0} 重跑的实验，目前计划中。

\section{后续计划}

\subsection{验证联合训练 vs 单模型}

当前的对比（EXP-04 vs EXP-05）因为 baseline 的梯度裁剪问题，结论不可靠。需要等 EXP-06（\texttt{grad\_clip=500.0} 的 baseline）跑完后，才能做公平对比。

核心问题：\textbf{logit 融合在 RL 中是否比单模型更有效？}

\subsection{梯度分析}

在跑实验的间隙，准备推导联合训练的梯度公式，验证其正确性。对于融合后的 policy $\pi_{\text{fused}}$，其 log-probability 为：
\begin{equation}
\log \pi_{\text{fused}}(a_t | s_t) = \log \text{softmax}\left((1-\lambda) \cdot z_1 + \lambda \cdot z_2\right)_{a_t}
\end{equation}

梯度关于 $\text{model}_2$ 参数 $\theta_2$ 的表达式：
\begin{equation}
\frac{\partial \log \pi_{\text{fused}}}{\partial \theta_2} = \lambda \left(\mathbf{e}_{a_t} - \pi_{\text{fused}}(\cdot|s_t)\right) \frac{\partial z_2}{\partial \theta_2}
\end{equation}

需要确认：这个梯度方向是否合理？$\lambda$ 的取值如何影响训练动态？

\subsection{探索更多 Weak-Driven 方法}

联合训练本质上是一种 \textbf{weak-driven learning}：让弱模型（$\text{model}_1$）在 logit 层提供参考信号，放大强模型（$\text{model}_2$）的学习效果。

这个思路可以推广到其他层次：

\begin{table}[h]
\centering
\begin{tabular}{lll}
\toprule
\textbf{层次} & \textbf{方法} & \textbf{弱模型的角色} \\
\midrule
Logit 层 & 当前方案（logit 融合） & 直接混合输出分布 \\
Reward / Sequence 层 & 弱模型作为 reward signal & 提供有指向性的优势估计 \\
Tool 层 & 弱模型参与工具调用 / 验证 & 在推理管道中提供辅助判断 \\
\bottomrule
\end{tabular}
\end{table}

核心原理是：弱模型提供的是\textbf{有结构的、有指向性的}梯度放大信号，而不是随机噪声。这种信号之所以有效，是因为即使弱模型本身不够强，它的知识分布与强模型是相关的——它们来自同一个预训练家族，共享底层的语言和推理模式。

\section{经验总结}

\begin{enumerate}[nosep]
    \item \textbf{GRPO 在联合训练中不 work}：EXP-01 跑了 100 步，GSM8K 0\% 准确率。原因可能是 GRPO 的 ratio-based 裁剪在融合 logits 下表现不稳定。
    \item \textbf{MiniRL 比 GRPO 更适合联合训练}：切换到 MiniRL 后立即收敛。MiniRL 的 REINFORCE 风格梯度（只通过 $\log \pi$）可能在 logit 融合场景下更稳定。
    \item \textbf{超参数不能直接迁移}：MiniRL 的 \texttt{seq-mean-token-sum} 聚合使梯度范数比 GRPO 大 100--500 倍，必须相应调整 \texttt{grad\_clip}。这是一个容易被忽略但影响巨大的问题。
    \item \textbf{离线评估很重要}：训练中的 validation（$n=1$）噪声大，尤其是在 AIME-2025 这样只有 30 个样本的基准上。搭建了离线 vLLM 推理 + pass@k/maj@k 评估管道后，结果更可靠。
    \item \textbf{Checkpoint 管理是刚需}：6 组实验消耗约 665 GB checkpoint 空间，必须有系统化的实验索引和清理策略。
\end{enumerate}

\begin{thebibliography}{9}
\bibitem{minirl}
Qwen Team.
\textit{Stabilizing Reinforcement Learning with LLMs: Formulation and Practices}.
arXiv:2512.01374, 2024.
\end{thebibliography}

\end{document}
